\Conclusion % заключение к отчёту

По результатам работы решены следующие задачи:
\begin{itemize}
  \item выполнена формальная постановка задачи трёхмерной паллетизации и введены метрические показатели: коэффициент заполнения (percolation), минимальный коэффициент опоры (min\_support\_ratio) и баланс высот (height\_balance);
  \item разработана и программно реализована вспомогательная структура данных HeightHandler для хранения карты высот с пятью различными реализациями (Baseline, Rects, Grid, SegTree, Quadtree);
  \item разработано пять алгоритмов укладки:
  \begin{itemize}
    \item GreedySolver --- жадный алгоритм с ограничением на устойчивость (56,04\% заполнения, 30,31\% опоры, 12,8~мс);
    \item DBLFSolver --- алгоритм Deepest Bottom Left Fill (65,19\% заполнения, 7,5~мс);
    \item FFDSolver --- алгоритм First Fit Decreasing (66,36\% заполнения, 8,3~мс);
    \item LNSSolver --- метаэвристический алгоритм на основе LNS с тринадцатью операторами мутации (63,49\% заполнения, 64,74\% опоры);
    \item GeneticSolver --- генетический алгоритм с популяционным подходом (63,45\% заполнения, 62,02\% опоры);
  \end{itemize}
  \item реализована поддержка мультипаллетного режима с метрикой баланса высот;
  \item проведено экспериментальное тестирование на 436 реальных задачах.
\end{itemize}

Экспериментальное сравнение пяти реализаций структуры HeightHandler показало, что простая реализация Rects на основе отсортированного списка прямоугольников превосходит более сложные структуры данных по производительности в 22 раза по сравнению с базовой реализацией.

Сравнительный анализ разработанных алгоритмов выявил две характерные группы. Детерминированные эвристики (DBLFSolver, FFDSolver) достигают наивысшего коэффициента заполнения (65--66\%) при времени выполнения менее 10~мс, однако не обеспечивают устойчивость размещения (минимальный коэффициент опоры $\approx$~2,4\%). Метаэвристические алгоритмы (LNSSolver, GeneticSolver) обеспечивают баланс между заполнением и устойчивостью. LNSSolver при лимите 5~с достигает коэффициента заполнения 63,49\% и минимального коэффициента опоры 64,74\%, превосходя GreedySolver на 7,45 и 34,43 процентных пункта соответственно. LNSSolver превосходит GeneticSolver на коротких временных интервалах благодаря оптимизации выборочного пересчёта паллет.

В мультипаллетном режиме LNSSolver демонстрирует перколяцию 62,09\% при балансе высот 84,26\%. FFDSolver показывает критически низкий баланс высот (52,16\%) из-за концентрации крупных коробок на первых паллетах.

Разработанное программное решение может быть интегрировано в WMS-систему склада для автоматической генерации схем паллетизации. Время выполнения алгоритма LNSSolver порядка 5 секунд на паллету является приемлемым для задач складской логистики.

В дальнейшем планируется исследовать адаптивный выбор весов операторов мутации (Adaptive LNS), обучение эвристики при помощи методов RL для online-сценариев, а также гибридный подход, комбинирующий высокую перколяцию FFDSolver с устойчивостью метаэвристических методов.
